% paper-to-book template (filled in by pipeline/p2b.py build)
\documentclass[11pt,letterpaper,openright]{book}
\usepackage[inner=1in,outer=2.3in,top=1in,bottom=1in,marginparwidth=1.75in,marginparsep=0.25in]{geometry}
\usepackage{amsmath,amssymb,amsthm,mathtools}
\usepackage{fontspec}
\usepackage{unicode-math}
\setmainfont{Latin Modern Roman}
\setmathfont{Latin Modern Math}
\newfontfamily\symbolfont{DejaVu Sans}\newcommand{\checkmk}{{\symbolfont ✓}}\newcommand{\crossmk}{{\symbolfont ✗}}
\usepackage{microtype}
\usepackage[dvipsnames]{xcolor}
\usepackage{longtable,booktabs,array,calc}
\usepackage{enumitem}
\usepackage{marginnote}
\usepackage{pdfpages}
\usepackage{tikz}
\usepackage{fancyhdr}
\usepackage{titlesec}
\usepackage[most]{tcolorbox}
\usepackage{imakeidx}
\makeindex[intoc,columns=2]
\usepackage[hidelinks]{hyperref}
\usepackage{bookmark}
\providecommand{\tightlist}{\setlength{\itemsep}{0pt}\setlength{\parskip}{0pt}}
\providecommand{\pandocbounded}[1]{#1}

\definecolor{unsourced}{RGB}{190,20,20}
\definecolor{srcgray}{RGB}{105,105,105}
\definecolor{leangreen}{RGB}{20,120,60}
\definecolor{quotebar}{RGB}{170,170,170}

% ---- margin tags: where each statement comes from ----
\newcommand{\srcmargin}[1]{\marginnote{\raggedright\scriptsize #1}}
\newcommand{\srcok}[3]{\textcolor{srcgray}{\textsf{Source:} \texttt{#1}\ifx&#2&\else, #2\fi}}
\newcommand{\srcbad}[1]{\textcolor{unsourced}{\textbf{\textsf{LEDGER CHECK FAILED}} (#1)}}
\newcommand{\unsourced}{\textcolor{unsourced}{\textbf{\textsf{UNSOURCED}}}}
\newcommand{\leanok}[1]{\textcolor{leangreen}{\textsf{Lean \checkmk} \texttt{\detokenize{#1}}}}
\newcommand{\leanbad}[1]{\textcolor{unsourced}{\textsf{Lean \crossmk} \texttt{\detokenize{#1}}}}
\newcommand{\leanpending}[1]{\textcolor{srcgray}{\textsf{not yet formally checked:} \texttt{\detokenize{#1}}}}
\newcommand{\srccite}[2]{\textcolor{srcgray}{[\texttt{#1}\ifx&#2&\else, #2\fi]}}
\newcommand{\derivedtag}[1]{\textcolor{srcgray}{\textsf{Derived here} from #1}}

% ---- the paper's HOL4 notation, so quotations from its LaTeX source print as in the paper ----
\providecommand{\HOLinline}[1]{\ensuremath{#1}}
\providecommand{\HOLFreeVar}[1]{\mathit{#1}}
\providecommand{\HOLBoundVar}[1]{\mathit{#1}}
\providecommand{\HOLConst}[1]{\textsf{#1}}
\providecommand{\HOLSymConst}[1]{#1}
\providecommand{\HOLNumLit}[1]{#1}
\providecommand{\HOLKeyword}[1]{\textbf{\textsf{#1}}}
\providecommand{\HOLFieldName}[1]{\textsf{#1}}
\providecommand{\HOLTokenExists}{\exists}\providecommand{\HOLTokenForall}{\forall}\providecommand{\HOLTokenUnique}{\exists!}
\providecommand{\HOLTokenConj}{\land}\providecommand{\HOLTokenDisj}{\lor}\providecommand{\HOLTokenNeg}{\lnot}
\providecommand{\HOLTokenImp}{\Rightarrow}\providecommand{\HOLTokenEquiv}{\Leftrightarrow}\providecommand{\HOLTokenLambda}{\lambda}
\providecommand{\HOLTokenIn}{\in}\providecommand{\HOLTokenNotIn}{\notin}\providecommand{\HOLTokenLt}{<}\providecommand{\HOLTokenLeq}{\le}
\providecommand{\HOLTokenNotEqual}{\neq}\providecommand{\HOLTokenCompose}{\circ}\providecommand{\HOLTokenTurnstile}{\vdash}
\providecommand{\HOLTokenDefEquality}{\mathrel{\stackrel{\scriptscriptstyle\mathrm{def}}{=}}}
\providecommand{\HOLTokenLeftbrace}{\{}\providecommand{\HOLTokenRightbrace}{\}}\providecommand{\HOLTokenBar}{\mid}
\providecommand{\HOLTokenEmpty}{\emptyset}

% ---- environments ----
\theoremstyle{plain}
\newtheorem{btheorem}{Theorem}[chapter]
\newtheorem{blemma}[btheorem]{Lemma}
\newtheorem{bproposition}[btheorem]{Proposition}
\newtheorem{bcorollary}[btheorem]{Corollary}
\theoremstyle{definition}
\newtheorem{bdefinitioninner}[btheorem]{Definition}
\newtheorem{bexample}[btheorem]{Example}
\newtheorem{bexercise}{Exercise}[chapter]
\theoremstyle{remark}
\newtheorem{bremark}[btheorem]{Remark}
\newenvironment{bdefinition}[1][]{\begin{tcolorbox}[enhanced,breakable,colback=black!3,colframe=black!3,
  boxrule=0pt,left=6pt,right=6pt,top=2pt,bottom=2pt,arc=0pt]\ifx&#1&\begin{bdefinitioninner}\else\begin{bdefinitioninner}[#1]\fi}
  {\end{bdefinitioninner}\end{tcolorbox}}
\newenvironment{bquote}{\begin{tcolorbox}[enhanced,breakable,colback=white,colframe=white,boxrule=0pt,
  borderline west={1.6pt}{0pt}{quotebar},left=10pt,right=4pt,top=1pt,bottom=1pt,arc=0pt,fontupper=\small\itshape]}
  {\end{tcolorbox}}
\newenvironment{binpaper}{\begin{tcolorbox}[enhanced,breakable,colback=blue!3,colframe=blue!25,boxrule=0.4pt,
  left=6pt,right=6pt,top=3pt,bottom=3pt,arc=1pt,title={\small\textsf{In the paper}},fonttitle=\color{black},
  colbacktitle=blue!8,coltitle=black]\small}{\end{tcolorbox}}

% ---- the paper's pages: note column on the right ----
\newcommand{\papernote}[1]{\begin{tikzpicture}[remember picture,overlay]
  \node[anchor=north west,text width=1.55in,font=\scriptsize\sffamily,align=left]
    at ([xshift=-2.0in,yshift=-1.1in]current page.north east) {\textbf{In this book}\\[3pt]#1};
  \end{tikzpicture}}

\pagestyle{fancy}
\fancyhf{}
\fancyhead[LE]{\small\scshape\leftmark}
\fancyhead[RO]{\small\scshape\rightmark}
\fancyfoot[LE,RO]{\small\thepage}
\renewcommand{\headrulewidth}{0pt}
\fancyhfoffset[LE,RO]{0pt}

\title{Windmills: A Prerequisite Book}
\begin{document}
\frontmatter
\begin{titlepage}\centering
\vspace*{2in}{\Huge Windmills: A Prerequisite Book\par}\vspace{0.4in}
{\Large A prerequisite book for Windmills of the minds: an algorithm for Fermat's Two Squares Theorem\par}\vspace{0.2in}{\large Hing Lun Chan\par}
\vfill{\small Assumed background: arithmetic (addition, subtraction, multiplication); high-school algebra; basic proof writing (induction, contradiction); sets and functions\par}
\vspace{0.3in}{\footnotesize Built with paper-to-book. Every definition and theorem is tagged in the margin
with its source; quotations are checked automatically against the retrieved sources, and
statements marked \textsf{Lean \checkmk} are verified against Lean's mathlib.\par}
\end{titlepage}
\tableofcontents
\chapter*{How to read this book}\addcontentsline{toc}{chapter}{How to read this book}
The chapters build, in order, everything the paper uses that goes beyond the assumed background.
Each definition and theorem carries a margin note saying where it comes from:
\begin{itemize}
\item \srcok{source}{§} --- quoted from a retrieved source, and the quotation was checked automatically;
\item \leanok{name} --- the statement matches a declaration in Lean's mathlib that compiles;
\item \leanpending{name} --- a formal counterpart exists but has not been checked yet;
\item \derivedtag{sources} --- not stated in any source, but proved in this book from the sourced results named, with the proof given and checked in Lean;
\item \unsourced{} or \srcbad{id} --- not backed by a checked source: treat with suspicion.
\end{itemize}
Shaded boxes are definitions; passages with a grey bar on the left are verbatim quotations;
blue boxes show where the idea appears in the paper. The last part is the paper itself.
\mainmatter
\hypertarget{whole-numbers}{%
\chapter{Whole Numbers}\label{whole-numbers}}

The paper proves a theorem of Fermat about primes and squares, and turns
the proof into an algorithm. In the paper's own words:

\begin{bquote}\srcmargin{\srcok{paper}{§1}{C1-16}}

Fermat's two squares theorem, dated back to 1640, states that a prime
\(n\) that is one more than a multiple of \(4\) can be uniquely
expressed as a sum of odd and even squares

\end{bquote}

To follow it you need five ideas about whole numbers: when one number
\emph{divides} another, \emph{division with remainder}, \emph{even and
odd}, \emph{primes}, and \emph{squares}. This chapter introduces each
and proves the few facts about them that the paper uses.

Throughout, a \emph{number} means a natural number
\(0, 1, 2, 3, \dots\), as in the paper:

\begin{bquote}\srcmargin{\srcok{paper}{§1.3}{C1-18}}

The set of natural numbers is denoted by \(\mathbb{N}\), counting from
\(0\)

\end{bquote}

Some of the definitions quoted below are stated for all integers,
negative ones included; we only ever apply them to natural numbers.

\hypertarget{divisibility}{%
\section{Divisibility}\label{divisibility}}

\begin{bdefinition}\label{def-divides}\srcmargin{\srcok{wp-divisor}{Definition}{L001}\\[2pt]\leanok{Dvd.dvd}}

\index{divisible}

A number \(m\) is \textbf{divisible} by a nonzero number \(d\) if
\(m = kd\) for some number \(k\). We then say \(d\) \emph{divides}
\(m\), and write \(d \mid m\).

\end{bdefinition}

\begin{bquote}\srcmargin{\srcok{wp-divisor}{Definition}{L001}}

An integer n is divisible by a nonzero integer m if there exists an
integer k such that n = k m .

\end{bquote}

So \(3 \mid 12\) because \(12 = 4 \cdot 3\), but \(3 \nmid 13\): no
number \(k\) has \(3k = 13\). Every number \(m \geq 1\) is divisible by
\(1\) and by itself, since \(m = m \cdot 1\). These two are its
\emph{trivial} divisors:

\begin{bquote}\srcmargin{\srcok{wp-divisor}{Terminology}{C1-01}}

1, −1, n and − n are known as the trivial divisors of n

\end{bquote}

A number greater than \(1\) with a divisor other than these is called
\emph{composite} \srccite{wp-divisor}{Terminology}.

\begin{bremark}

The paper writes \(n \mid m\) for ``\(n\) divides \(m\)'', with \(n\) as
the divisor, and from here on we use its letters (§1.3):

\begin{bquote}\srcmargin{\srcok{paper}{§1.3}{C1-17}}

We write
\HOLinline{\HOLFreeVar{n}\;\HOLConst{\ensuremath{\mid}}\;\HOLFreeVar{m}}
when \(n\) divides \(m\), which is equivalent to
\HOLinline{\HOLFreeVar{m}\;\ensuremath{\equiv}\;\HOLNumLit{0}\;\ensuremath{(}\ensuremath{\bmod}\;\HOLFreeVar{n}\ensuremath{)}}
when
\HOLinline{\HOLFreeVar{n}\;\HOLSymConst{\HOLTokenNotEqual{}}\;\HOLNumLit{0}}.

\end{bquote}

Lean's mathlib, which checks the statements in this book, defines
\(d \mid m\) for every \(d\), including \(d = 0\): there \(0 \mid m\)
holds only when \(m = 0\). For \(d \ne 0\) the two definitions agree.

\end{bremark}

\hypertarget{division-with-remainder}{%
\section{Division with remainder}\label{division-with-remainder}}

Most numbers don't divide each other exactly; what is left over is the
\emph{remainder}.

\begin{btheorem}\label{thm-division}\srcmargin{\srcok{wp-euclidean-division}{Division theorem}{L002}\\[2pt]\leanok{Book.Windmills.Ch01.division}\\\leanok{Book.Windmills.Ch01.division\_unique}}

\textbf{Division with remainder.} Given numbers \(m\) and \(n\) with
\(n \neq 0\), there are unique numbers \(q\) (the \emph{quotient}) and
\(r\) (the \emph{remainder}) such that
\[m = nq + r \qquad\text{and}\qquad 0 \le r < n.\]

\end{btheorem}

\begin{bquote}\srcmargin{\srcok{wp-euclidean-division}{Division theorem}{L002}}

Given two integers a and b , with b ≠ 0 , there exist unique integers q
and r such that a = b q + r , and 0 ≤ r \textless{} \textbar{} b
\textbar{}

\end{bquote}

\begin{bexample}

Dividing \(17\) by \(3\): the quotient is \(5\) and the remainder is
\(2\), since \(17 = 3 \cdot 5 + 2\) and \(2 < 3\).
\srccite{wp-euclidean-division}{Intuitive example}

\end{bexample}

\begin{proof}

\emph{Existence.} Start with \(r = m\) and \(q = 0\), so that
\(m = nq + r\). While \(r \ge n\), replace \(r\) by \(r - n\) and \(q\)
by \(q + 1\); the equation \(m = nq + r\) still holds. Each step makes
\(r\) smaller, and a number can only get smaller finitely many times, so
the process stops, with \(r < n\).

\emph{Uniqueness.} Suppose also \(m = nq' + r'\) with \(r' < n\), and
say \(r \le r'\) (otherwise swap the two). Then \(nq = nq' + (r' - r)\),
so \(q \ge q'\) and \(r' - r = n(q - q')\). The left side is less than
\(n\), and the only multiple of \(n\) less than \(n\) is \(0\). So
\(r' = r\), and then \(q' = q\).

\end{proof}

The paper writes the quotient as \(m \operatorname{div} n\) and the
remainder as \(m \operatorname{mod} n\):

\begin{bquote}\srcmargin{\srcok{paper}{§1.3}{C1-05}}

The integer quotient and remainder of \(m\) divided by \(n\) are written
as \HOLinline{\HOLFreeVar{m}\;\HOLConst{\HOLConst{div}}\;\HOLFreeVar{n}}
and
\HOLinline{\HOLFreeVar{m}\;\HOLConst{\HOLConst{mod}}\;\HOLFreeVar{n}},
respectively.

\end{bquote}

For example, \(17 \operatorname{div} 3 = 5\) and
\(17 \operatorname{mod} 3 = 2\). Dividing leaves nothing over exactly
when \(n\) divides \(m\):

\begin{bproposition}\label{prop-dvd-mod}\srcmargin{\derivedtag{\texttt{wp-divisor}, \texttt{wp-euclidean-division}}\\[2pt]\leanok{Book.Windmills.Ch01.dvd\_iff\_mod}}

For \(n \ne 0\): \(\ n \mid m\) if and only if
\(m \operatorname{mod} n = 0\).

\end{bproposition}

\begin{proof}

If \(m = kn\), then \(m = nk + 0\) with \(0 < n\), so by uniqueness the
remainder is \(0\). Conversely, if the remainder is \(0\), then
\(m = nq\), so \(n \mid m\).

\end{proof}

\hypertarget{even-and-odd}{%
\section{Even and odd}\label{even-and-odd}}

\begin{bdefinition}\label{def-parity}\srcmargin{\srcok{wp-parity}{Definition}{C1-06}\\\srcok{wp-parity}{lead}{C1-20}\\[2pt]\leanok{Even}\\\leanok{Odd}}

\index{even}

A number is \textbf{even} if it has the form \(2k\), and \textbf{odd} if
it has the form \(2k + 1\), for some number \(k\). Whether a number is
even or odd is called its \textbf{parity}.

\end{bdefinition}

\begin{bquote}\srcmargin{\srcok{wp-parity}{Definition}{C1-06}}

An even number is an integer of the form x = 2 k where k is an integer;
{[} 4 {]} an odd number is an integer of the form x = 2 k + 1

\end{bquote}

\begin{bquote}\srcmargin{\srcok{wp-parity}{lead}{C1-20}}

parity is the property of an integer of whether it is even or odd

\end{bquote}

\begin{btheorem}\label{thm-even-odd}\srcmargin{\derivedtag{\texttt{wp-euclidean-division}, \texttt{wp-parity}}\\[2pt]\leanok{Book.Windmills.Ch01.even\_or\_odd}\\\leanok{Book.Windmills.Ch01.not\_even\_and\_odd}\\\leanok{Book.Windmills.Ch01.even\_iff\_mod}\\\leanok{Book.Windmills.Ch01.odd\_iff\_mod}}

Every number is either even or odd, and not both. A number \(n\) is even
exactly when \(n \operatorname{mod} 2 = 0\), and odd exactly when
\(n \operatorname{mod} 2 = 1\).

\end{btheorem}

\begin{proof}

Divide \(n\) by \(2\): \(n = 2q + r\) with \(r \in \{0, 1\}\). If
\(r = 0\) then \(n = 2q\) is even; if \(r = 1\) then \(n = 2q + 1\) is
odd. So every number is one or the other.

Conversely, if \(n = 2k\) then \(n = 2k + 0\), and if \(n = 2k + 1\)
then that is already a division by \(2\) with remainder \(1\). By
uniqueness of the remainder, an even number has remainder \(0\) and an
odd number has remainder \(1\); in particular no number is both.

\end{proof}

Parity behaves simply under addition: even \(+\) even and odd \(+\) odd
are even, while even \(+\) odd is odd (write the numbers as \(2a\) or
\(2a + 1\) and add). Parity matters throughout the paper, both of
numbers and, from Chapter 5 on, of the \emph{sizes of sets}:

\begin{bquote}\srcmargin{\srcok{paper}{Abstract}{C1-19}}

The theory behind involves involutions on a finite set, especially the
parity of the number of fixed points in the involutions.

\end{bquote}

\begin{bremark}

mathlib phrases these definitions slightly differently: \texttt{Even\ n}
says \(n = r + r\) for some \(r\), and \texttt{Odd\ n} says
\(n = 2k + 1\). The first is the same as \(n = 2r\).

\end{bremark}

\begin{binpaper}

The paper's notation section (§1.3) writes \(\mathsf{even}\ n\) and
\(\mathsf{odd}\ n\) for these two properties:

\begin{bquote}\srcmargin{\srcok{paper}{§1.3}{C1-08}}

prime n means it is a prime, and even n or odd n denotes its parity

\end{bquote}

\end{binpaper}

\hypertarget{prime-numbers}{%
\section{Prime numbers}\label{prime-numbers}}

\begin{bdefinition}\label{def-prime}\srcmargin{\srcok{wp-prime}{lead}{L005}\\[2pt]\leanok{Nat.Prime}}

\index{prime number}

A \textbf{prime number} is a number greater than \(1\) that is not a
product of two smaller numbers.

\end{bdefinition}

\begin{bquote}\srcmargin{\srcok{wp-prime}{lead}{L005}}

A prime number (or a prime ) is a natural number greater than 1 that is
not a product of two smaller natural numbers.

\end{bquote}

The \emph{divisors} of \(n\) are the numbers that divide it evenly
\srccite{wp-prime}{Definition and examples}. Primes are exactly the
numbers greater than \(1\) with no divisors except the trivial ones:

\begin{bproposition}\label{prop-prime-divisors}\srcmargin{\srcok{wp-prime}{Definition and examples}{C1-09}\\[2pt]\leanok{Nat.prime\_def}}

A number \(p\) is prime if and only if it has exactly two positive
divisors, \(1\) and \(p\) (so in particular \(p \ge 2\)).

\end{bproposition}

\begin{proof}

First, \(0\) and \(1\) are not prime (the definition needs \(p > 1\)),
and neither has exactly the two positive divisors \(1\) and \(p\): every
positive number divides \(0\), and \(1\) has the single positive divisor
\(1\). So take \(p > 1\).

If \(p\) is a product \(ab\) of two numbers smaller than \(p\), then
\(a \ne 1\) (else \(b = p\), which is not smaller) and \(a \ne p\), so
\(a\) is a positive divisor other than \(1\) and \(p\).

Conversely, suppose \(p\) has a positive divisor \(a\) other than \(1\)
and \(p\), say \(p = ab\). Then \(b \ge 1\) (as \(p \ne 0\)), so
\(a \le p\); since \(a \ne p\) and \(a \ne 1\), we have \(1 < a < p\).
Then \(b = p / a\) satisfies \(1 < b < p\) as well, so \(p\) is a
product of two smaller numbers.

\end{proof}

In the same way, a number greater than \(1\) is composite exactly when
it is not prime.

\begin{bexample}

\(7\) is prime: none of \(2, 3, 4, 5, 6\) divides it
\srccite{wp-prime}{Definition and examples}. The primes below \(30\) are
\(2, 3, 5, 7, 11, 13, 17, 19, 23, 29\).

\end{bexample}

Two facts about primes matter for the paper.

\begin{btheorem}\label{thm-only-even-prime}\srcmargin{\srcok{paper}{§2}{C1-12}\\[2pt]\leanok{Book.Windmills.Ch01.prime\_even\_eq\_two}}

The only even prime is \(2\).

\end{btheorem}

\begin{bquote}\srcmargin{\srcok{paper}{§2}{C1-12}}

The only even prime is
\HOLinline{\HOLNumLit{2}\;\HOLSymConst{=}\;\HOLNumLit{1}\ensuremath{{\sp{\HOLNumLit{2}}}}\;\HOLSymConst{\ensuremath{+}}\;\HOLNumLit{1}\ensuremath{{\sp{\HOLNumLit{2}}}}},
a sum of two squares.

\end{bquote}

\begin{proof}

An even number \(n\) has \(2 \mid n\). If \(n\) is prime, its only
divisors are \(1\) and \(n\), so \(n = 2\).

\end{proof}

\begin{btheorem}\label{thm-odd-prime-mod-4}\srcmargin{\srcok{paper}{§2}{C1-13}\\[2pt]\leanok{Book.Windmills.Ch01.odd\_prime\_mod\_four}}

An odd prime leaves a remainder of \(1\) or \(3\) when divided by \(4\).

\end{btheorem}

\begin{proof}

Divide the prime \(p\) by \(4\): \(p = 4q + r\) with
\(r \in \{0, 1, 2, 3\}\). If \(r\) were \(0\) or \(2\), then
\(p = 2(2q)\) or \(p = 2(2q + 1)\) would be even. Since \(p\) is odd,
\(r\) is \(1\) or \(3\).

\end{proof}

So the odd primes fall into two families: those of the form \(4k + 1\)
(\(5, 13, 17, 29, \dots\)) and those of the form \(4k + 3\)
(\(3, 7, 11, 19, 23, \dots\)). The paper's theorem is about the first
family.

\begin{binpaper}

§2 opens with exactly these two facts: the only even prime is a sum of
two squares, and the odd primes split by their remainder on division by
\(4\).

\end{binpaper}

\hypertarget{squares}{%
\section{Squares}\label{squares}}

\begin{bdefinition}\label{def-square}\srcmargin{\srcok{wp-square-number}{lead}{L006}\\[2pt]\leanok{IsSquare}}

\index{square number}

A \textbf{square number} (or perfect square) is the square
\(k^2 = k \cdot k\) of a number \(k\).

\end{bdefinition}

\begin{bquote}\srcmargin{\srcok{wp-square-number}{lead}{L006}}

a square number or perfect square is an integer that is the square of an
integer

\end{bquote}

The first few are \(0, 1, 4, 9, 16, 25, 36, 49, 64\). A \emph{sum of two
squares} is a number of the form \(u^2 + v^2\); for instance
\(13 = 3^2 + 2^2\) and \(29 = 5^2 + 2^2\), while no choice of \(u, v\)
gives \(3\), \(7\) or \(11\) (the squares up to \(11\) are only
\(0, 1, 4, 9\), and no two of them add up to \(3\), \(7\) or \(11\)).
Chapter 2 explains why.

\begin{binpaper}

The paper writes \(\mathsf{square}\ n\) for ``\(n\) is a square''
(§1.3):

\begin{bquote}\srcmargin{\srcok{paper}{§1.3}{C1-15}}

For a natural number \(n \in \mathbb{N}\),
\HOLinline{\HOLConst{square}\;\HOLFreeVar{n}} means it is a square:
\HOLinline{\HOLSymConst{\HOLTokenExists{}}\HOLBoundVar{k}.\;\HOLFreeVar{n}\;\HOLSymConst{=}\;\HOLBoundVar{k}\HOLSymConst{\ensuremath{\sp{2}}}}

\end{bquote}

\end{binpaper}

\hypertarget{exercises}{%
\section{Exercises}\label{exercises}}

\begin{bexercise}

Find the quotient and remainder when \(29\) is divided by \(4\), and
when \(61\) is divided by \(4\). Both are prime; which family of odd
primes does each belong to?

\end{bexercise}

\begin{bexercise}

Show that if \(n \mid a\) and \(n \mid b\) then \(n \mid a + b\).

\end{bexercise}

\begin{bexercise}

Show that the square of an odd number is odd. (Write the number as
\(2k + 1\).)

\end{bexercise}

\begin{bexercise}

Write each of \(5, 13, 17, 29, 37, 41\) as a sum of two squares.

\end{bexercise}

\part{The Paper}
\chapter*{Reader's guide to the paper}\addcontentsline{toc}{chapter}{Reader's guide to the paper}
The paper follows exactly as published. In the margin of each page are the concepts from this book that appear on it, with the chapter that teaches them. The first appearance of each:

\begin{longtable}{p{0.5\textwidth}ll}
Divisibility & Chapter~1 & p.~1 of the paper\\
Even and odd numbers & Chapter~1 & p.~1 of the paper\\
Prime numbers & Chapter~1 & p.~1 of the paper\\
Division with remainder & Chapter~1 & p.~2 of the paper\\
Square numbers & Chapter~1 & p.~2 of the paper\\
\end{longtable}
\includepdf[pages=1,scale=0.86,offset=-0.55in 0in,pagecommand={\papernote{Divisibility \hfill Ch.~1\\Even and odd numbers \hfill Ch.~1\\Prime numbers \hfill Ch.~1}}]{/home/claude/paper-to-book/books/windmills/sources/paper/paper.pdf}
\includepdf[pages=2,scale=0.86,offset=-0.55in 0in,pagecommand={\papernote{Division with remainder \hfill Ch.~1\\Even and odd numbers \hfill Ch.~1\\Prime numbers \hfill Ch.~1\\Square numbers \hfill Ch.~1}}]{/home/claude/paper-to-book/books/windmills/sources/paper/paper.pdf}
\includepdf[pages=3,scale=0.86,offset=-0.55in 0in,pagecommand={\papernote{Division with remainder \hfill Ch.~1\\Even and odd numbers \hfill Ch.~1\\Prime numbers \hfill Ch.~1\\Square numbers \hfill Ch.~1}}]{/home/claude/paper-to-book/books/windmills/sources/paper/paper.pdf}
\includepdf[pages=4,scale=0.86,offset=-0.55in 0in,pagecommand={\papernote{Prime numbers \hfill Ch.~1\\Square numbers \hfill Ch.~1}}]{/home/claude/paper-to-book/books/windmills/sources/paper/paper.pdf}
\includepdf[pages=5,scale=0.86,offset=-0.55in 0in,pagecommand={\papernote{\textit{}}}]{/home/claude/paper-to-book/books/windmills/sources/paper/paper.pdf}
\includepdf[pages=6,scale=0.86,offset=-0.55in 0in,pagecommand={\papernote{\textit{}}}]{/home/claude/paper-to-book/books/windmills/sources/paper/paper.pdf}
\includepdf[pages=7,scale=0.86,offset=-0.55in 0in,pagecommand={\papernote{Division with remainder \hfill Ch.~1\\Even and odd numbers \hfill Ch.~1\\Prime numbers \hfill Ch.~1\\Square numbers \hfill Ch.~1}}]{/home/claude/paper-to-book/books/windmills/sources/paper/paper.pdf}
\includepdf[pages=8,scale=0.86,offset=-0.55in 0in,pagecommand={\papernote{\textit{}}}]{/home/claude/paper-to-book/books/windmills/sources/paper/paper.pdf}
\includepdf[pages=9,scale=0.86,offset=-0.55in 0in,pagecommand={\papernote{Divisibility \hfill Ch.~1\\Division with remainder \hfill Ch.~1\\Even and odd numbers \hfill Ch.~1}}]{/home/claude/paper-to-book/books/windmills/sources/paper/paper.pdf}
\includepdf[pages=10,scale=0.86,offset=-0.55in 0in,pagecommand={\papernote{Divisibility \hfill Ch.~1\\Division with remainder \hfill Ch.~1\\Even and odd numbers \hfill Ch.~1}}]{/home/claude/paper-to-book/books/windmills/sources/paper/paper.pdf}
\includepdf[pages=11,scale=0.86,offset=-0.55in 0in,pagecommand={\papernote{Division with remainder \hfill Ch.~1\\Even and odd numbers \hfill Ch.~1\\Prime numbers \hfill Ch.~1}}]{/home/claude/paper-to-book/books/windmills/sources/paper/paper.pdf}
\includepdf[pages=12,scale=0.86,offset=-0.55in 0in,pagecommand={\papernote{Division with remainder \hfill Ch.~1\\Even and odd numbers \hfill Ch.~1\\Prime numbers \hfill Ch.~1\\Square numbers \hfill Ch.~1}}]{/home/claude/paper-to-book/books/windmills/sources/paper/paper.pdf}
\includepdf[pages=13,scale=0.86,offset=-0.55in 0in,pagecommand={\papernote{Prime numbers \hfill Ch.~1}}]{/home/claude/paper-to-book/books/windmills/sources/paper/paper.pdf}
\includepdf[pages=14,scale=0.86,offset=-0.55in 0in,pagecommand={\papernote{Division with remainder \hfill Ch.~1\\Prime numbers \hfill Ch.~1\\Square numbers \hfill Ch.~1}}]{/home/claude/paper-to-book/books/windmills/sources/paper/paper.pdf}
\backmatter
\chapter{Sources}
Every source used, as retrieved for this book. The identifiers in the margins refer to this list.
\begin{description}\item[\texttt{ml-function-basic}] https://leanprover-community.github.io/mathlib4\_docs/Mathlib/Logic/Function/Basic.html. \url{https://leanprover-community.github.io/mathlib4_docs/Mathlib/Logic/Function/Basic.html} --- 1 of 1 quotations verified (retrieved 2026-11-12).
\item[\texttt{ml-periodic}] https://leanprover-community.github.io/mathlib4\_docs/Mathlib/Dynamics/PeriodicPts/Defs.html. \url{https://leanprover-community.github.io/mathlib4_docs/Mathlib/Dynamics/PeriodicPts/Defs.html} --- 3 of 3 quotations verified (retrieved 2026-11-12).
\item[\texttt{paper}] Windmills of the minds: an algorithm for Fermat's Two Squares Theorem. \url{arXiv:2112.02556} --- 48 of 48 quotations verified (retrieved 2026-11-12).
\item[\texttt{wp-bijection}] https://en.wikipedia.org/wiki/Bijection. \url{https://en.wikipedia.org/wiki/Bijection} --- 1 of 1 quotations verified (retrieved 2026-11-12).
\item[\texttt{wp-cardinality}] https://en.wikipedia.org/wiki/Cardinality. \url{https://en.wikipedia.org/wiki/Cardinality} --- 2 of 2 quotations verified (retrieved 2026-11-12).
\item[\texttt{wp-cartesian}] https://en.wikipedia.org/wiki/Cartesian\_product. \url{https://en.wikipedia.org/wiki/Cartesian_product} --- 1 of 1 quotations verified (retrieved 2026-11-12).
\item[\texttt{wp-composition}] https://en.wikipedia.org/wiki/Function\_composition. \url{https://en.wikipedia.org/wiki/Function_composition} --- 1 of 1 quotations verified (retrieved 2026-11-12).
\item[\texttt{wp-divisor}] https://en.wikipedia.org/wiki/Divisor. \url{https://en.wikipedia.org/wiki/Divisor} --- 3 of 3 quotations verified (retrieved 2026-11-12).
\item[\texttt{wp-euclidean-division}] https://en.wikipedia.org/wiki/Euclidean\_division. \url{https://en.wikipedia.org/wiki/Euclidean_division} --- 3 of 3 quotations verified (retrieved 2026-11-12).
\item[\texttt{wp-fermat-two-squares}] https://en.wikipedia.org/wiki/Fermat\%27s\_theorem\_on\_sums\_of\_two\_squares. \url{https://en.wikipedia.org/wiki/Fermat%27s_theorem_on_sums_of_two_squares} --- 4 of 4 quotations verified (retrieved 2026-11-12).
\item[\texttt{wp-finite-set}] https://en.wikipedia.org/wiki/Finite\_set. \url{https://en.wikipedia.org/wiki/Finite_set} --- 1 of 1 quotations verified (retrieved 2026-11-12).
\item[\texttt{wp-fixed-point}] https://en.wikipedia.org/wiki/Fixed\_point\_(mathematics). \url{https://en.wikipedia.org/wiki/Fixed_point_(mathematics)} --- 1 of 1 quotations verified (retrieved 2026-11-12).
\item[\texttt{wp-hol}] https://en.wikipedia.org/wiki/HOL\_(proof\_assistant). \url{https://en.wikipedia.org/wiki/HOL_(proof_assistant)} --- 2 of 2 quotations verified (retrieved 2026-11-12).
\item[\texttt{wp-inverse}] https://en.wikipedia.org/wiki/Inverse\_function. \url{https://en.wikipedia.org/wiki/Inverse_function} --- 2 of 2 quotations verified (retrieved 2026-11-12).
\item[\texttt{wp-involution}] https://en.wikipedia.org/wiki/Involution\_(mathematics). \url{https://en.wikipedia.org/wiki/Involution_(mathematics)} --- 9 of 9 quotations verified (retrieved 2026-11-12).
\item[\texttt{wp-iterated}] https://en.wikipedia.org/wiki/Iterated\_function. \url{https://en.wikipedia.org/wiki/Iterated_function} --- 2 of 2 quotations verified (retrieved 2026-11-12).
\item[\texttt{wp-lambda}] https://en.wikipedia.org/wiki/Lambda\_calculus. \url{https://en.wikipedia.org/wiki/Lambda_calculus} --- 1 of 1 quotations verified (retrieved 2026-11-12).
\item[\texttt{wp-loop-invariant}] https://en.wikipedia.org/wiki/Loop\_invariant. \url{https://en.wikipedia.org/wiki/Loop_invariant} --- 2 of 2 quotations verified (retrieved 2026-11-12).
\item[\texttt{wp-modular}] https://en.wikipedia.org/wiki/Modular\_arithmetic. \url{https://en.wikipedia.org/wiki/Modular_arithmetic} --- 3 of 3 quotations verified (retrieved 2026-11-12).
\item[\texttt{wp-orbit}] https://en.wikipedia.org/wiki/Orbit\_(dynamics). \url{https://en.wikipedia.org/wiki/Orbit_(dynamics)} --- 1 of 1 quotations verified (retrieved 2026-11-12).
\item[\texttt{wp-parity}] https://en.wikipedia.org/wiki/Parity\_(mathematics). \url{https://en.wikipedia.org/wiki/Parity_(mathematics)} --- 5 of 5 quotations verified (retrieved 2026-11-12).
\item[\texttt{wp-periodic-point}] https://en.wikipedia.org/wiki/Periodic\_point. \url{https://en.wikipedia.org/wiki/Periodic_point} --- 3 of 3 quotations verified (retrieved 2026-11-12).
\item[\texttt{wp-permutation}] https://en.wikipedia.org/wiki/Permutation. \url{https://en.wikipedia.org/wiki/Permutation} --- 2 of 2 quotations verified (retrieved 2026-11-12).
\item[\texttt{wp-prime}] https://en.wikipedia.org/wiki/Prime\_number. \url{https://en.wikipedia.org/wiki/Prime_number} --- 4 of 4 quotations verified (retrieved 2026-11-12).
\item[\texttt{wp-proof-assistant}] https://en.wikipedia.org/wiki/Proof\_assistant. \url{https://en.wikipedia.org/wiki/Proof_assistant} --- 1 of 1 quotations verified (retrieved 2026-11-12).
\item[\texttt{wp-set-builder}] https://en.wikipedia.org/wiki/Set-builder\_notation. \url{https://en.wikipedia.org/wiki/Set-builder_notation} --- 1 of 1 quotations verified (retrieved 2026-11-12).
\item[\texttt{wp-square-number}] https://en.wikipedia.org/wiki/Square\_number. \url{https://en.wikipedia.org/wiki/Square_number} --- 1 of 1 quotations verified (retrieved 2026-11-12).
\item[\texttt{wp-turnstile}] https://en.wikipedia.org/wiki/Turnstile\_(symbol). \url{https://en.wikipedia.org/wiki/Turnstile_(symbol)} --- 1 of 1 quotations verified (retrieved 2026-11-12).
\item[\texttt{wp-unique}] https://en.wikipedia.org/wiki/Uniqueness\_quantification. \url{https://en.wikipedia.org/wiki/Uniqueness_quantification} --- 1 of 1 quotations verified (retrieved 2026-11-12).
\item[\texttt{wp-while}] https://en.wikipedia.org/wiki/While\_loop. \url{https://en.wikipedia.org/wiki/While_loop} --- 1 of 1 quotations verified (retrieved 2026-11-12).\end{description}
\printindex
\end{document}
